\documentclass[11pt]{article}

\usepackage[margin=1in]{geometry}
\usepackage{amsmath}
\usepackage{amssymb}
\usepackage{booktabs}
\usepackage{graphicx}
\usepackage{microtype}
\usepackage{placeins}
\usepackage{xcolor}
\usepackage[hidelinks]{hyperref}

\newcommand{\fOne}{\(f_1\)}
\newcommand{\fThree}{\(f_3\)}
\newcommand{\KL}{D_{\mathrm{KL}}}
\newcommand{\resultgraphic}[1]{%
  \IfFileExists{#1}{%
    \includegraphics[width=\linewidth]{#1}%
  }{%
    \fbox{\parbox[c][0.18\textheight][c]{0.92\linewidth}{%
      \centering Missing generated figure: \texttt{#1}}}%
  }%
}
\newcommand{\resulttable}[1]{%
  \InputIfFileExists{#1}{}{%
    \textit{Missing generated table: \texttt{#1}}%
  }%
}

\title{Leakage-Controlled Punctuation Stylometry for Human Authorship,\\
LLM Imitation, and Detection}
\author{Authors omitted for review}
\date{September 2026}

\begin{document}
\maketitle

\begin{abstract}
We evaluate punctuation as a deliberately narrow instrument for measuring
human author identifiability, prompted literary-style imitation, and
human--LLM separation.  The primary analysis uses 20 authors with three
content-distinct books each, 400 Gemini~2.5 Flash generations, and 200
Gemini~2.5 Pro generations.  It fixes adjacent-punctuation features
(\fThree), 2,000 punctuation marks, leave-one-book-out references, author-equal
estimands, and author-blocked uncertainty before inference.  Human passages
are attributed to the correct author with 73.1\% macro accuracy (95\%
author-bootstrap interval 59.3--85.0\%) against 5\% chance.  Prompted
generations are attributed to their requested author only 8.75\% for Flash and
8.0\% for Pro; their paired difference is unresolved.  Nevertheless,
out-of-author detection separates generated from human passages well
(AUC 0.952 for Flash and 0.931 for Pro), with true-positive rates of 75.5\%
and 71.5\% at an empirical human false-positive rate of 5.36\%.  A
three-book cross-fit also links a human-only separability diagnostic to
held-out attribution (\(\rho=0.53\)).  Finally, target divergence increases
from the first to the fifth 1,000-mark window for both Flash and Pro.  These
results show that detectable non-human punctuation and successful imitation
of a requested author's punctuation are distinct properties.  The
contribution is the integrated, leakage-controlled evaluation instrument;
neither KL divergence, the punctuation features, nor the language models are
new.
\end{abstract}

\section{Introduction}

Punctuation is a low-dimensional trace of writing style.  It carries
information about rhythm, clause structure, dialogue, and formatting while
discarding almost all lexical content.  Earlier work showed that sequences of
punctuation marks support human authorship analysis
\cite{darmon2021,stamatatos2009}.  Large language models create a related but
different test: a generated passage can be recognisably non-human without
matching the punctuation habits of the author named in its prompt.

We study three questions.  First, can held-out human passages be assigned to
their authors using punctuation alone?  Second, when a model receives an
author excerpt as a style reference, does its output enter the target author's
punctuation region?  Third, can deviation from a target-specific human
reference distinguish generated from genuine passages for authors that did
not calibrate the detector?  We also ask whether a human-only separability
score predicts prospective attribution and whether imitation changes with
position in a generation.

The analysis was motivated by a frozen ten-author campaign.  That campaign is
retained as discovery and replication evidence, not pooled into the present
headline estimands.  Its Flash drift estimate was unresolved, whereas the
balanced 20-author analysis detects increasing target KL for both Flash and
Pro.  The primary claims below come from the declared 20-author, cross-fitted
inference package.

The main findings are:
\begin{itemize}
  \item Human punctuation is strongly author-specific in this screened panel:
  author-equal 20-way attribution is 73.1\% at 2,000 marks.
  \item Prompted imitation is weak in absolute target-attribution terms:
  Flash and Pro score 8.75\% and 8.0\%, respectively, against 5\% chance, with
  wide author-level intervals and no resolved model contrast.
  \item Weak target imitation coexists with strong human--LLM detection.  At a
  threshold calibrated on other authors, AUC exceeds 0.93 for both models.
  \item Cross-fitted \(D/C\) is positively associated with held-out human
  attribution, while target KL rises across generation position for both
  models.
\end{itemize}

\paragraph{Novelty boundary.}
We do not introduce the punctuation representation, KL divergence,
likelihood-ratio testing, resampling, or the evaluated models.  Our
contribution is their combination into a declared evaluation instrument that
keeps source books out of their own references, calibrates detection on other
authors, makes authors the population unit, validates screening prospectively,
and separates exploratory ten-author evidence from primary 20-author
estimands.  This boundary matters because leakage control and the choice of
population unit change the substantive conclusions.

\section{Related work}

Authorship attribution commonly combines lexical, character, syntactic, and
structural cues; Stamatatos surveys the broader field
\cite{stamatatos2009}.  Darmon et al.\ isolate punctuation sequences and define
the frequency and transition-derived features used here
\cite{darmon2021}.  Our nearest-reference attribution is intentionally simpler
than modern high-dimensional authorship systems: its purpose is to audit a
specific stylistic channel.

Generated-text detection is sensitive to the generator, domain, decoding
scheme, and access assumptions.  Prior work has examined human judgements and
automatic features \cite{ippolito2020}, generator attribution
\cite{uchendu2021}, perturbation-based zero-shot detection
\cite{mitchell2023}, and fundamental limits under distribution shift and
paraphrase \cite{sadasivan2024}.  We therefore report a conditional detector,
not a universal one.  Its threshold is learned empirically from genuine human
texts and tested on held-out authors; no asymptotic punctuation null is assumed.

\section{Study design}
\label{sec:design}

\subsection{Discovery analysis and primary population}

The frozen discovery campaign contained ten authors and an unbalanced
22-book corpus.  It exposed the importance of leave-one-book-out references
and motivated chunk-size, feature, and drift analyses.  Those outputs remain a
historical benchmark.  They do not define the current headline population and
are not used to select among the primary results.

The primary panel contains 20 English-language prose authors and exactly three
content-distinct books per author.  It combines the ten discovery authors with
ten additional authors screened using a full-corpus human-only \fThree\
\(D/C\) diagnostic.  Before LLM evaluation, each proposed author had to
achieve at least two correct whole-book attributions out of three under a
20-way leave-one-book-out check; all ten passed.  This yields 60 books and
5\% attribution chance.  The screen creates a conditional evaluation
population enriched for punctuation separability.  It is useful for testing
the instrument but is not a random sample of authors.

\subsection{Generated texts}

Each generation is conditioned on a 1,500-word excerpt from one of the target
author's source books and asks for an original continuation in that style.
The campaign uses temperature 1.0 and continues generation until at least
5,000 analysed punctuation marks are available.  There are 20 Flash runs and
10 Pro runs per author, for 400 and 200 runs in total.  The first ten authors'
runs are reused from the frozen campaign and the ten extension authors were
generated later.  Thus model-level results use all 20 authors, but
old-versus-new cohort comparisons would mix author composition with run date
or provider state and are not causal.

For continuity with the discovery campaign, analysed text replaces em and en
dashes with commas.  Exact raw responses were retained only for the extension
cohort.  There, the transformation changes comma share by 0.0083 percentage
points for Flash and 0.0013 percentage points for Pro, so it is too small to
explain the reported extension-cohort results.  The absence of raw historical
responses prevents a full-panel audit of this policy.

\subsection{Declared primary and sensitivity specifications}

The declared primary specification is the full 20-author panel, \fThree,
2,000 punctuation marks, leave-one-book-out references, add-0.5 smoothing for
detection, and a 5\% human false-positive target.  Accuracy estimands weight
authors equally.  Sensitivities vary feature (\fOne\ versus \fThree), chunk
size (1,000, 2,000, and 4,000), author cohort, and smoothing (0.1, 0.5, and
1.0).  Pooled-reference calculations are retained only as leakage contrasts.
Run-level binomial tests are descriptive and are not treated as
population-level inference.

\section{Methods}
\label{sec:methods}

\subsection{Punctuation representation and references}

Text is reduced to the ten-mark alphabet
\[
  \{!,\,"\, ,(\,),\, ,\,.\, ,:\, ;\, ,?\, ,\hat{\ }\},
\]
where the caret is the parser's end marker.  The \fOne\ vector contains
single-mark frequencies.  The primary \fThree\ vector is the flattened joint
distribution of adjacent marks, obtained by scaling each row of the
first-order transition matrix by the marginal frequency of its initial mark
\cite{darmon2021}.  Attribution assigns an observation \(x\) to the reference
author \(a\) minimizing directed divergence,
\[
  \widehat a(x)=\arg\min_a \KL\!\left(f_3(x)\,\|\,f_3(R_a)\right).
\]

For a human passage, the complete source book is excluded from its target
author's reference, leaving the other two books.  For an LLM run, the book
that supplied the prompt excerpt is excluded from the requested author's
reference.  Exclusion is resolved by content identity, and each scored
observation is checked against the reference audit.  This prevents a source
book from helping to recognise itself through its remaining text.

\subsection{Attribution estimands and uncertainty}

Human books are divided into non-overlapping chunks; incomplete tails are
dropped.  Generated-text attribution uses the first \(m\) marks of each run
for chunk size \(m\).  Within each source type and model we first calculate an
accuracy for every author and then average those rates, so prolific books and
the larger Flash run count do not change author weights.  The primary
confidence intervals are central 95\% percentile intervals from 2,000
author-block bootstrap draws.  Resampling an author carries all of that
author's books, chunks, prompt-source groups, and runs together.  Per-author
diagnostic intervals additionally resample books for human text and prompt
source books for LLM text.

The Flash--Pro contrast is paired by author.  We report the author-equal mean
of Flash minus Pro target-attribution accuracy, a paired author bootstrap
interval, and a two-sided Wilcoxon signed-rank test.  This contrast does not
mistake 600 repeated generations for 600 independent population units.

\subsection{Out-of-author detection}

Detection uses a target-specific likelihood-ratio \(G\) score.  For \fThree,
observed transition counts are compared row by row with an add-\(\epsilon\)
smoothed reference transition profile.  Larger scores indicate greater
departure from the requested author's human reference.  The nominal
\(\chi^2\) approximation is a diagnostic only because adjacent punctuation
transitions and longer-range prose structure violate independent-trial
assumptions.

Authors are allocated to five folds, stratified by original versus extension
cohort.  For each fold, the operating threshold is the 95th percentile of
human \(G\) scores from the other four folds.  That threshold is then applied
only to human and generated passages from held-out authors.  Fold-specific
scores are divided by their calibration thresholds before pooling the
out-of-fold ROC curve.  The bootstrap resamples authors and re-estimates every
fold threshold on each valid draw.  We report out-of-fold AUC, empirical human
false-positive rate, and LLM true-positive rate.

\subsection{Prospective separability}

The three books per author support a separate three-way cross-fit.  On each
fold, two books define the author profile and within-author inconsistency
\[
  C_a=\tfrac12\{\KL(f_{a1}\|f_{a2})+\KL(f_{a2}\|f_{a1})\}.
\]
The distance \(D_a\) is the directed KL from that two-book profile to its
nearest competing author profile.  Neither \(D_a\), \(C_a\), nor the
references use the held-out third book.  We geometrically average \(D_a/C_a\)
over the three held-out-book rotations, attribute the excluded books and their
2,000-mark chunks, and calculate the across-author Spearman association
between the ratio and chunk accuracy.  A 2,000-draw author bootstrap supplies
the interval for \(\rho\).  This is a prospective validation of a screening
diagnostic, although the panel itself remains selected.

\subsection{Positional drift}

Every run contributes five consecutive, non-overlapping 1,000-mark windows.
For each model, feature, author, and run we subtract the first-window metric
from the fifth-window metric and then average runs within author.  Author-equal
means and author-block bootstrap intervals summarize the change.  Directional
Wilcoxon signed-rank tests evaluate increasing target KL and rank and
decreasing target margin and hit rate; Benjamini--Hochberg adjustment is
applied jointly to the 16 model--feature--metric tests
\cite{benjamini1995}.  Position is not randomized, so ``drift'' denotes an
association with generation position rather than a causal context-length
effect.

\section{Results}
\label{sec:results}

\subsection{Primary attribution}

\begin{table}[htbp]
  \centering
  \resizebox{\textwidth}{!}{%
    \resulttable{tables/inference_v2/primary_attribution.tex}}
  \caption{Primary 20-author, \fThree, 2,000-mark attribution.  Accuracy is
  author-equal; intervals resample authors.  Human units are held-out chunks,
  while LLM units are independent runs.  Chance is 5\%.}
  \label{tab:primary-attribution}
\end{table}

Table~\ref{tab:primary-attribution} establishes a large gap between human
identifiability and successful prompted imitation.  Human macro accuracy is
73.1\% (95\% CI 59.3--85.0\%), whereas Flash reaches 8.75\%
(3.25--15.75\%) and Pro 8.0\% (1.0--19.5\%).  The latter intervals include
chance and should not be converted into claims that either model reliably
imitates the population of requested authors.  The paired Flash-minus-Pro
difference is 0.75 percentage points (95\% CI \(-11.76\) to 11.25; signed-rank
\(p=0.395\)), which resolves neither superiority nor equivalence.

\begin{figure}[htbp]
  \centering
  \resultgraphic{figures/inference_v2/attribution_by_author.pdf}
  \caption{Per-author target-attribution accuracy for human passages, Flash,
  and Pro under the primary specification.  The dashed line is 20-way chance.
  Heterogeneity across authors motivates author-level uncertainty and cautions
  against interpreting the pooled run count as the inferential sample size.}
  \label{fig:attribution}
\end{figure}

Figure~\ref{fig:attribution} shows that the macro averages conceal substantial
heterogeneity.  Generated passages also collapse toward a few predicted
authors: in the primary analysis, Wells is the nearest reference for 148 of
400 Flash runs and 115 of 200 Pro runs.  This concentration is evidence of a
shared model-side geometry, not successful imitation of those predicted
authors.

\subsection{Chunk-size sensitivity}

\begin{table}[htbp]
  \centering
  \begin{tabular}{lrrr}
    \toprule
    Source & 1,000 marks & 2,000 marks (primary) & 4,000 marks \\
    \midrule
    Human & 65.8 & 73.1 & 73.6\(^\dagger\) \\
    Flash & 12.5 & 8.75 & 8.5 \\
    Pro & 11.5 & 8.0 & 7.0 \\
    \bottomrule
  \end{tabular}
  \caption{Author-equal \fThree\ attribution accuracy (\%) by chunk size.
  \(^\dagger\)Only 19 authors supply complete 4,000-mark human chunks, so this
  endpoint is not a perfectly matched full-panel comparison.}
  \label{tab:length-attribution}
\end{table}

Human attribution rises from 65.8\% at 1,000 marks to 73.1\% at 2,000 and is
similar at 4,000 marks, subject to the missing short-corpus author.  LLM target
attribution does not improve with length: Flash changes from 12.5\% to 8.75\%
to 8.5\%, and Pro from 11.5\% to 8.0\% to 7.0\%.  Thus 2,000 marks is the
declared primary compromise rather than a post hoc optimum.

\subsection{Detection and empirical calibration}

\begin{table}[htbp]
  \centering
  \resizebox{\textwidth}{!}{%
    \resulttable{tables/inference_v2/detection.tex}}
  \caption{Five-fold out-of-author detection for the primary \fThree,
  2,000-mark specification.  Each held-out author is scored at a threshold
  calibrated only on other authors.  Intervals use an author-block bootstrap
  that re-estimates fold thresholds.}
  \label{tab:detection}
\end{table}

The detector distinguishes generated from human punctuation even though
target-author imitation is near chance.  Flash has AUC 0.952 (95\% CI
0.891--0.984) and Pro 0.931 (0.857--0.974).  At the declared 5\% operating
point, the held-out human FPR is 5.36\% for both conditions, while TPR is
75.5\% for Flash and 71.5\% for Pro.  The author-bootstrap TPR intervals are
wide (30.0--95.8\% and 27.0--93.5\%), indicating that detector performance is
not uniform across authors.

\begin{figure}[htbp]
  \centering
  \resultgraphic{figures/inference_v2/detection_oof_roc.pdf}
  \caption{Out-of-author ROC curves.  Scores from each held-out fold are
  normalized by the human threshold learned from the other folds; dots mark
  the declared 5\% operating point.}
  \label{fig:roc}
\end{figure}

Length sensitivity is directionally consistent.  For Flash, AUC/TPR is
0.903/52.8\% at 1,000 marks and 0.970/87.5\% at 4,000 marks; for Pro it is
0.852/39.0\% and 0.962/87.5\%.  The corresponding empirical human FPRs are
4.99\% and 7.11\%, respectively.  The 4,000-mark improvement is therefore
accompanied by less precise tail calibration and fewer human chunks.  At
2,000 marks, changing the smoothing constant from 0.1 to 1.0 changes Flash
AUC only from 0.951 to 0.953 and Pro AUC from 0.929 to 0.932.

\begin{figure}[htbp]
  \centering
  \resultgraphic{figures/inference_v2/g_calibration.pdf}
  \caption{Failure of nominal \(\chi^2\) calibration for the \(G\) statistic
  versus empirical out-of-author calibration at 2,000 marks.  Under
  leave-one-book-out human references, nominal 5\% rejection is 95.3\% for
  \fOne\ and 71.1\% for \fThree; the cross-fitted empirical procedure returns
  the held-out rate to approximately 5\%.}
  \label{fig:g-calibration}
\end{figure}

Figure~\ref{fig:g-calibration} explains why cross-fitting is essential.
Dependent punctuation and book-to-book heterogeneity make the nominal
likelihood-ratio reference severely anti-conservative.  A detector calibrated
against that reference would report impressive but invalid rejection rates.

\subsection{Cross-fitted human separability}

\begin{table}[htbp]
  \centering
  \resizebox{\textwidth}{!}{%
    \resulttable{tables/inference_v2/separability.tex}}
  \caption{Three-book cross-fitted separability.  The primary row is \fThree;
  \fOne\ is a declared feature sensitivity.  Accuracy is author-equal.}
  \label{tab:separability}
\end{table}

For \fThree, the geometric-mean \(D/C\) score has Spearman
\(\rho=0.534\) with held-out chunk attribution (95\% author-bootstrap CI
0.021--0.895; ordinary Spearman \(p=0.015\)).  Held-out whole-book macro
accuracy is 80.0\% and chunk macro accuracy is 72.0\%.  The \fOne\ sensitivity
is similar (\(\rho=0.576\), interval 0.082--0.936).  These results support
\(D/C\) as a useful ranking diagnostic, not as a universal threshold or an
unbiased estimate of performance on unscreened authors.

\begin{figure}[htbp]
  \centering
  \resultgraphic{figures/inference_v2/crossfit_separability.pdf}
  \caption{Human-only cross-fitted \(D/C\) against held-out 2,000-mark
  attribution.  Each point is one author; the excluded book contributes to
  neither coordinate's reference construction.}
  \label{fig:separability}
\end{figure}

\subsection{Positional drift in both models}

\begin{table}[htbp]
  \centering
  \resizebox{\textwidth}{!}{%
    \resulttable{tables/inference_v2/drift.tex}}
  \caption{Author-equal change from the first to fifth 1,000-mark window.
  Intervals resample authors; \(p\)-values are directional signed-rank tests
  adjusted across the declared drift family.}
  \label{tab:drift}
\end{table}

Target KL increases by 0.092 for Flash (95\% CI 0.046--0.144,
BH-adjusted \(p=0.00143\)) and by 0.168 for Pro (0.125--0.214,
adjusted \(p=3.05\times10^{-5}\)).  The corresponding target-hit changes,
\(-4.5\) and \(-7.5\) percentage points, have intervals that include zero and
adjusted \(p\)-values of 0.111 and 0.0725.  The continuous divergence endpoint
therefore detects positional movement more clearly than the coarse 20-way hit
indicator.  Crucially, the current result is drift in \emph{both} Flash and
Pro; the unresolved discovery estimate does not describe this panel.

\begin{figure}[htbp]
  \centering
  \resultgraphic{figures/inference_v2/positional_drift.pdf}
  \caption{Author-equal target KL across five consecutive 1,000-mark windows.
  Ribbons are author-block bootstrap intervals.  Both models move away from
  requested-author punctuation profiles with generation position.}
  \label{fig:drift}
\end{figure}

\FloatBarrier
\section{Discussion}

\subsection{Imitation and detection are different estimands}

The apparent tension between near-chance target attribution and high detection
AUC is informative.  Attribution asks whether a generated passage is closer
to its requested author than to 19 competitors.  Detection asks whether the
passage has an unusually large target-relative \(G\) score compared with
genuine held-out human writing.  A model can fail to reproduce the requested
punctuation geometry while remaining systematically outside the human
target-reference distribution.  Indeed, concentration of generated
predictions on a few authors and increasing target KL are consistent with a
model-side punctuation tendency overwhelming author-specific conditioning.

More marks improve human attribution and especially detection, but do not
improve requested-author attribution.  That pattern argues against describing
the issue as simple sampling noise.  It is instead consistent with weak or
decaying author-specific control in this feature space.  The positional design
cannot isolate the mechanism: later windows also occur later in the narrative,
after more self-generated context, and nearer a possible ending.

\subsection{What the evaluation instrument adds}

Four design choices prevent common overstatements.  First, source books are
excluded from their own target references by content identity.  Second,
detection thresholds are trained on other authors and evaluated out of fold.
Third, authors rather than chunks or runs determine population uncertainty.
Fourth, the full panel, feature, length, smoothing, and operating point are
declared, while pooled-reference and alternate-grid outputs are labelled as
sensitivities or leakage contrasts.  Together these choices turn a collection
of descriptive distances into estimands with explicit target populations and
uncertainty.

\section{Limitations}

\begin{enumerate}
  \item \textbf{Selected population.}  Half of the panel was screened for high
  full-corpus \fThree\ \(D/C\) and all new authors passed an attribution gate.
  Headline estimates are conditional on this 20-author panel and likely
  overstate performance for a broad, unscreened author population.

  \item \textbf{Corpus scope.}  The corpus contains public-domain English prose
  with only three books per author.  Period, genre, editorial practice,
  dialogue share, and transcription conventions may be confounded with
  authorship.  Chunk counts vary sharply by author even though macro
  estimands restore equal author weight.

  \item \textbf{Few population clusters.}  Twenty authors support cluster-aware
  inference, but the wide intervals correctly reflect limited information at
  the population level.  At 4,000 marks one author has no complete human
  chunk, so that sensitivity is not a matched 20-author endpoint.

  \item \textbf{Generator and prompt scope.}  The results cover Gemini~2.5
  Flash and Pro under one excerpt-conditioned generation protocol and one
  temperature.  They do not establish behavior for other models, prompts,
  languages, or decoding settings.  The ten-author and extension generations
  were not contemporaneous.

  \item \textbf{Preprocessing.}  Dash-to-comma accounting is exact only for the
  extension cohort because raw frozen responses are unavailable.  Although
  the observed perturbation is tiny, a complete full-panel sensitivity cannot
  be reconstructed.

  \item \textbf{Internal rather than external validation.}  Author folds block
  direct calibration leakage, but all folds come from the same constructed
  corpus and campaign.  A genuinely prospective study should freeze the
  instrument and evaluate new authors, books, providers, and generation dates.

  \item \textbf{Non-causal drift.}  Window position jointly changes generated
  context, narrative phase, and proximity to termination.  Randomized stopping
  lengths or controlled continuation experiments are needed to identify a
  causal context-length effect.

  \item \textbf{Narrow stylistic channel.}  Punctuation cannot establish
  authorship, literary quality, memorization, or ``human-ness'' in general.
  The detector may be evaded by editing or paraphrase
  \cite{sadasivan2024}, and its operating characteristics should not be
  exported beyond the evaluated domain.
\end{enumerate}

\section{Conclusion}

In a declared 20-author evaluation, punctuation alone identifies held-out
human authors well, but Gemini~2.5 Flash and Pro rarely reproduce the requested
author's punctuation closely enough for 20-way attribution.  Those same
outputs are readily separated from genuine target-author writing when
thresholds are calibrated on other authors.  Cross-fitted \(D/C\) predicts
where human attribution will work, and target divergence increases with
generation position for both models.  The central methodological lesson is
that source exclusion, out-of-author calibration, author-level resampling, and
a fixed claim hierarchy are not implementation details: they determine what
the experiment can validly say.

\begin{thebibliography}{99}

\bibitem{benjamini1995}
Y.~Benjamini and Y.~Hochberg.
\newblock Controlling the false discovery rate: a practical and powerful
approach to multiple testing.
\newblock \emph{Journal of the Royal Statistical Society: Series B},
57(1):289--300, 1995.

\bibitem{darmon2021}
A.~Darmon, M.~Bazzi, S.~D. Howison, and M.~A. Porter.
\newblock Pull out all the stops: textual analysis via punctuation sequences.
\newblock \emph{European Journal of Applied Mathematics}, 32(6):1069--1105,
2021.

\bibitem{ippolito2020}
D.~Ippolito, D.~Duckworth, C.~Callison-Burch, and D.~Eck.
\newblock Automatic detection of generated text is easiest when humans are
fooled.
\newblock In \emph{Proceedings of the 58th Annual Meeting of the Association
for Computational Linguistics}, pages 1808--1822, 2020.

\bibitem{mitchell2023}
E.~Mitchell, Y.~Lee, A.~Khadilkar, C.~D. Manning, and C.~Finn.
\newblock DetectGPT: zero-shot machine-generated text detection using
probability curvature.
\newblock In \emph{Proceedings of the 40th International Conference on Machine
Learning}, pages 24950--24962, 2023.

\bibitem{sadasivan2024}
V.~S. Sadasivan, A.~Kumar, S.~Balasubramanian, W.~Wang, and S.~Feizi.
\newblock Can AI-generated text be reliably detected?
\newblock \emph{Transactions on Machine Learning Research}, 2024.

\bibitem{stamatatos2009}
E.~Stamatatos.
\newblock A survey of modern authorship attribution methods.
\newblock \emph{Journal of the American Society for Information Science and
Technology}, 60(3):538--556, 2009.

\bibitem{uchendu2021}
A.~Uchendu, Z.~Ma, T.~Le, R.~Zhang, and D.~Lee.
\newblock Authorship attribution for neural text generation.
\newblock In \emph{Proceedings of the 2021 Conference on Empirical Methods in
Natural Language Processing}, pages 8384--8395, 2021.

\end{thebibliography}

\end{document}
